\documentstyle[% 


righttag% 


,11pt% 


,amssymb% 


,russian%               remove in Eng. 


,izvru%                 Set izven% in Eng. 


]{amsart} 


 


% verify missed $$ line 337


\newcommand{\la}{\langle} 


\newcommand{\ra}{\rangle} 


\newcommand{\tts}[1]{}  


 


 


%\hoffset=-15mm%                  For Eng. 


\setcounter{page}{67} 


\god{2003} 


\nomer{8~(495)} %                        Remove in Eng. 


%\nomer{Vol.\,47, No.\,8}%               Set in Eng. 


%\str{\pageref{begin}--\pageref{end}}%  Set in Eng. 


\udk{517.956} 


 


\author{\it Е.М.\,Богатов, О.М.\,Пенкин} 


% NB:   ^^ remove in Eng. 


\title{Метод Фурье в задаче Коши для 


уравнения четвертого порядка на стратифицированных множествах} 


\thanks{Работа выполнена при финансовой поддержке Российского фонда 


фундаментальных исследований (проект 01-01-00417).} 


 


\date{} 


%\pagestyle{myheadings}%         Set in Eng. 


\pagestyle{plain} %              Remove in Eng. 


\begin{document} 


%\thispagestyle{plain}%          Set in Eng. 


\maketitle 


\label{begin} 


 


 


В последние годы наметился устойчивый интерес к описанию статических и 


динамических явлений в системах составного типа (в механическом 


контексте речь может идти, например, о колебаниях конструкции, 


составленной из струн и мембран). Простейший пример такого типа был 


pассмотрен Р.\,Курантом \cite{1} еще в 30-е годы, систематические же 


исследования в этом направлении начались сравнительно недавно. 


 


Основная трудность в изучении упомянутых систем состоит в том, что с 


геометрической точки зрения они не являются многообразиями. Существуют 


два способа обойти эту трудность. Один из них принадлежит математикам 


из Западной Европы (G.\,Lumer, S.\,Nicaise, F.\,Ali Mehmeti, J.\,von Below)


и подробно описан в \cite{2}. Второй способ --- в случае, когда система 


геометрически представляет собой граф, --- был разработан 


Ю.В.\,Покорным \cite{3}, \cite{4}, в многомерном варианте ---


авторами \cite{5}--\cite{10}. 


 


В данной статье для уравнения 4-го порядка на множествах, составленных 


из конечного числа многообразий (стратифицированных множествах), 


обосновывается метод Фурье. Нам удобно придерживаться второго из 


упомянутых выше подходов, основанного на введении на стратифицированном 


множестве аналога классической дивергенции. Ради автономности статьи 


остановимся на некоторых деталях этого подхода. 


\medskip 


 


{\it Вспомогательные сведения}. 


Как и в цитированных выше работах, 


страты, составляющие 


стратифицированное множество 


$\Omega \subset {\bold R}^n$, считаются гладкими подмногообразиями 


${\bold R}^n$ и обозначаются $\sigma_{ki}$ ($k$ --- размерность, 


$i$ --- номер). Примыкание страта $\sigma_{ki}$ к 


$\sigma_{mj}$ ($m>k$) записывается в виде 


$\sigma_{ki}\prec\sigma_{mj}$. Далее считается также, что замыкание каждого 


$\sigma_{ki}$ компактно. 


 


 


Через $\Omega_0$ обозначим открытое связное (в топологии $\Omega$, 


индуцированной из ${\bold R}^n$) подмножество $\Omega$, составленное 


из его стратов, замыкание которого совпадает с $\Omega$. 


Тогда множество $\partial\Omega_0=\Omega\setminus\Omega_0$ 


является границей $\Omega_0$ в указанной топологии. 


 


Каждый страт $\sigma_{ki}$ ($k\geq 0$) может рассматриваться как 


риманово многообразие с метрикой $(g_{\alpha\beta})$, индуцированной 


на $\sigma_{ki}$ из ${\bold R}^n$. Определим оператор дивергенции на 


гладких касательных векторных полях $\Vec F$ на $\Omega_0$ (т.\,е. 


$\Vec F(x)$ принадлежит касательному пространству $T_x\sigma_{ki}$ при 


$x\in\sigma_{ki}$ для любых $k$, $i$) соотношением 


\begin{equation} 


(\nabla\vec F)(x)= 


\frac{1}{\sqrt{g}}\frac{\partial}{\partial y^\alpha} 


(\sqrt{g}g^{\alpha\beta}F_\beta)(x)+ 


\sum_{\sigma_{kj}\succ\sigma_{k-1i}} 


\frac{1}{\sqrt{g^{kk}}} 


g^{k\alpha}F_\alpha\Big|_{\overline{kj}}(x). 


\end{equation} 


Здесь $x\in\sigma_{k-1i}$, $F_\alpha$ --- ковариантные компоненты 


$\Vec F$. Локальные координаты вблизи $x$ 


определяются только на $\sigma_{k-1i}$ и на всех 


$\sigma_{kj}\succ\sigma_{k-1i}$, причем делается это отдельно для 


каждой пары $(\sigma_{k-1i};\sigma_{kj})$ как на многообразии с 


краем ($\sigma_{k-1i}$ --- часть края $\overline\sigma_{kj}$), 


так что $y^k\equiv 0$ на $\sigma_{k-1i}$ и $y^k>0$ на 


$\sigma_{kj}$. 


 


Далее через 


$u\big|_{\overline{kj}}$ всегда обозначается продолжение по 


непрерывности на $\overline\sigma_{kj}$ сужения $u$ на $\sigma_{kj}$. 


Первое слагаемое в (1) совпадает с классической $(k-1)$-мерной 


дивергенцией $\nabla_{k-1}\Vec F$ в точке $x\in\sigma_{k-1i}$ (0-мерный 


классический оператор дивергенции считаем нулевым), а второе представляет 


собой сумму скалярных произведений $(\Vec F\cdot\Vec\nu)\big|_{kj}(x)$, 


где $\Vec\nu$ --- единичная нормаль в точке $x$, направленная внутрь 


$\sigma_{kj}$. Таким образом, формулу (1) можно переписать в виде 


$$ 


(\nabla\Vec F)(x)=(\nabla_{k-1}\Vec F)(x)+ 


\sum_{\sigma_{kj}\succ\sigma_{k-1i}}(\Vec F\cdot\Vec\nu) 


\Big|_{\ov{kj}}(x). 


$$ 


Определим класс $C^m_\sigma(\Omega_0)$ функций $u: 


\Omega_0\rightarrow{\bold R}$, 


$m$ раз непрерывно дифференцируемых на 


каждом страте, как на многообразии, причем их 


производные $(m-1)$-го порядка допускают 


продолжение по непрерывности в страты 


$\sigma_{k-1i}\prec\sigma_{kj}$, не входящие в $\partial\Omega_0$, 


и аналогичный ему класс функций $C^m_\sigma(\Omega)$. 


Скажем теперь, что касательное векторное поле гладко, если его 


компоненты принадлежат множеству $C^1_\sigma (\Omega_0)$. 


 


Обозначив через $\nabla u$ градиент $u$, заметим, что при $u\in 


C^2_{\sigma}(\Omega_0)$ определен оператор 


$\Delta u=\nabla(\nabla u)$, который является аналогом обычного 


оператора Лапласа--Бельтрами. 


Наряду с $\Delta$ будем рассматривать оператор $\Delta_p=\nabla(p\nabla)$ 


при $p\in C^1_\sigma(\Omega_0)$, т.\,е. 


$$ 


(\Delta_p u)(x)=\nabla_{k-1}(p\nabla u)(x)+(p\nabla u)_{\Vec\nu}(x), 


$$ 


где 


\begin{equation} 


(p\nabla u)_{\Vec\nu}(x)= 


\sum_{\sigma_{kj}\succ\sigma_{k-1i}}\bigg(p 


\frac{\partial u}{\partial\nu}\Big|_{\overline{kj}}\bigg)(x) 


\end{equation} 


--- аналог нормальной производной. 


 


Если сужение функции $f:\Omega \to {\bold R}$  


на каждый страт является равномерно 


непрерывной функцией (не являясь, вообще говоря, непрерывной в целом на 


$\Omega$), то будем называть еe постратно непрерывной.  


Множество таких функций 


обозначим $C_\sigma(\Omega)$. Для функций из $C_\sigma(\Omega)$ 


определим интеграл 


$$


\int_{\Omega_0}f\,d\mu=\sum_{\sigma_{ki}\subset\Omega_0} 


\ \int_{\sigma_{ki}}f\,d\mu, 


$$


где в правой части стоят обычные интегралы Римана по стандартной мере на 


каждом страте. 


Формулой 


$\la f,g\ra_\mu=\int\limits_{\Omega_0}fg\,d\mu$


задается скалярное произведение на $C_\sigma(\Omega)$. 


Определяемая им норма обозначается через $\|\cdot\|_\mu$. Теперь 


можем ввести новое пространство 


$L^2_\mu(\Omega_0)$ --- 


пополнение $C_\sigma(\Omega)$ по норме $\|\cdot\|_\mu$. 


 


Введем еще несколько функциональных пространств: 


\begin{itemize} 


\item[] $C(\Omega_0)$ --- класс функций $u: 


\Omega_0\to {\bold R}$, непрерывных на $\Omega_0$, и 


аналогичный ему класс функций $C(\Omega)$;


\item[] 


$C^m(\Omega_0)=C^m_\sigma(\Omega_0)\cap C(\Omega_0)$ и аналогичный 


ему класс функций $C^m(\Omega)$. 


\end{itemize} 


\medskip 


 


 


{\it О спектре дифференциального оператора $4$-го порядка ${\cal L}$ 


на $\Omega$ $($стационарный случай$)$.} 


Рассмотрим на $C^4(\Omega)$ дифференциальный оператор 


$$ 


{\cal L}u=\Delta_{p}(r\Delta_{p})u-\Delta_{q}u+k u, 


$$ 


$k(x) \in L^2_{\mu}(\Omega_0)$, $k(x)\geq 0$; $r \in \overline 


C{}^2_\sigma(\Omega_0)$, $q\in C^1_\sigma(\Omega_0)$ 


cтрого положительны; а также задачу 


\begin{alignat}{2} 


{\cal L}u(x)&=\lambda u(x), &\quad  x &\in \Omega_0, \\ 


(p\nabla u)_{\Vec\nu}(x)&=0, &\quad  x &\in\Omega\setminus\sigma_k, \\ 


u(x)=(p\nabla u)_{\Vec\nu}(x)&=0, &\quad x &\in \partial\Omega_0, 


\end{alignat} 


где в (4) стоит сумма, аналогичная (2), однако под еe знаком 


отсутствуют страты, входящие в $\partial\Omega_0$. 


Решением краевой задачи вида (3)--(5) в случае, когда $\Omega\subset\bold 


R^2$, является амплитудная функция свободных колебаний натянутой 


пластинчато-стержневой системы\footnote{Постановка задачи на модельном 


примере имеется в (\cite{11}, гл.\,3).}, помещенной на упругое основание. 


Условие (4) в этой ситуации означает, что стержни и пластины жестко 


сочленены друг с другом, поэтому будем называть его {\it условием 


жесткости}. В общем случае оно запрещает поворачиваться друг относительно 


друга элементам сопряжения стратифицированной конструкции, как это могло 


бы происходить при шарнирном соединении. 


 


Наличие члена $\Delta_q u$ в уравнении (3), связанного с натяжением 


системы, или члена $ku$, показывающего наличие упругого основания, обычно 


обусловлено тем, что при их отсутствии система может оказаться статически 


неопределенной, а задача (3)--(5) --- вырожденной, как показано в работах 


\cite{3}, \cite{12}. Правда, в нашем случае имеется еще условие жесткости, 


которое сохраняет упомянутую статическую определимость при $k\equiv 


q\equiv0$. Но даже перечисленных трех условий недостаточно для однозначной 


разрешимости задачи (3)--(5); требуется еще условие так называемой 


прочности множества $\Omega$.





Назовем $\Omega$ {\it прочным}, если любой 


страт $\sigma_{ki}\subset\Omega_0$ можно соединить с некоторым стратом 


$\sigma_{mj}\subset\partial\Omega_0$ связной цепочкой стратов 


$\sigma_{k_1m_1},\dotsc,\sigma_{k_p m_p}$, обладающих следующими свойствами: 


\begin{itemize} 


\item[] $\sigma_{k_1m_1}=\sigma_{ki}$, $\sigma_{k_p m_p}=\sigma_{mj}$; 


\item[] $\sigma_{k_q m_q}\subset\Omega_0$ для любого $q$, $1<q<p$; 


\item[] для любого $q$, $1\le q\le p$, либо $\sigma_{k_q 


i_q}\succ\sigma_{k_{q+1}i_{q+1}}$, либо $\sigma_{k_q i_q}\prec 


\sigma_{k_{q+1}i_{q+1}}$; 


\item[] $|k_{q+1}-k_q|=1$ для любого $q$, $1\le q\le p$. 


\end{itemize} 


Подобного рода условие имеется в работах \cite{5}, \cite{6}. В 


приведенной форме это условие сформулировано в \cite{13}, всюду далее 


оно предполагается выполненным. 


 


Следующая лемма (формула Грина для оператора $\cal L$) доказана на основе 


результатов, полученных в \cite{14}. 


 


\begin{lem} 


Пусть $w\in C^4_\sigma(\Omega)$, $v\in C^2_{rg}(\Omega_0)$


--- пространство функций из 


$C^2(\Omega)$, удовлетворяющих условию жесткости $(4)$.


Тогда 


\begin{multline} 


\int_{\Omega_0}v\cal L w\,d\mu=\int_{\Omega_0} 


(k v w+q\nabla v\nabla w+r\Delta_p v\Delta_p w)d\mu+ \\ 


% 


+\int_{\partial\Omega_0}v(q\nabla w)_{\Vec\nu}d\mu- 


\int_{\partial\Omega_0}[v(p\nabla(r\Delta_p w))_{\Vec\nu} 


-r\Delta_p w(p\nabla v)_{\Vec \nu}]d\mu. 


\end{multline} 


\end{lem} 


 


Собственным значением ${\cal L}$, соответствующим собственному 


элементу $u \in \overset{\circ}{H}{}^2(\Omega_0)$, называется число 


$\lambda$, удовлетворяющее для любого $\eta \in 


\overset{\circ}{H}{}^2(\Omega_0)$ равенству 


$\la{\cal L}u,\eta\ra_\mu=\lambda\la u,\eta\ra_\mu$, 


индуцированному формулой (6). 


 


Пространство $\overset{\circ}{H}{}^2(\Omega_0)$ получается 


пополнением $\overset{\circ}{C}{}^2_{rg}(\Omega_0)$ (подмножества 


$C^2_{rg}(\Omega_0)$), образованного финитными на $\Omega_0$ 


функциями, по норме $\|u\|_2=\sqrt{(u,u)}$, 


порождeнной скалярным произведением 


$$


(u,v)=\int_{\Omega_0}{(r\Delta_p 


u\Delta_p v+q\nabla u \nabla v + k u v)}d\mu. 


$$


 


\begin{thm} 


Спектральная задача $(3)$--$(5)$ имеет счетное 


множество решений $\lambda=\lambda_k$, $u=u_{k}(x)$, $k=1,2,\dotsc$, в 


пространстве $\overset{\circ}{H}{}^2(\Omega_0)$, каждое из которых 


вещественно и имеет конечную кратность. Собственные значения 


$\lambda_k$ можно расположить в порядке их возрастания${:}$ 


$0<\lambda_1\leq\lambda_2\leq\dotsb\le\lambda_n\leq\dotsb$, $\lambda_n 


\to \infty$. 


Cобственные функции $\{u_{k}(x)\}$ 


образуют базис в $L^2_{\mu}(\Omega_0)$ и $\overset{\circ}{H}{}^2(\Omega_0)$, 


ортонормированный в $L^2_{\mu}(\Omega_0)$ и ортогональный в смысле 


скалярного произведения $(\ ,\ )$. 


\end{thm} 


 


{\it Динамика\/${:}$ постановка задачи.} 


Определим $Q_T$ как стратифицированный цилиндр, 


состоящий из стратов $\sigma_{ij}^T=\sigma_{ij}\times (0,T)$. 


Основной интерес представляет нахождение 


функции $u(x,t)$ в цилиндре $Q_T= \Omega \times (0,T)$, 


удовлетворяющей требованиям 


\begin{gather} 


{\cal L}u+ u_{tt}=f(x,t),\\ 


(p\nabla u)_{\Vec\nu}(x,t)=0,\quad 


(x,t)\in Q_T\setminus\sigma^T_k, \\ 


u(x,0)=\varphi(x),\quad u_t(x,0)=\psi(x),\\ 


u|_{S_T}=(p\nabla u)_{\Vec\nu}|_{S_T}=0, 


\end{gather} 


где $S_T=\partial \Omega_0 \times [0,T]$, а в (8) стоит сумма, аналогичная 


сумме для средней компоненты равенства из (10), однако под ее знаком 


отсутствуют страты, входящие в $S_T$. Будем считать, что $\psi(x)\in 


L^2_\mu(\Omega_0)$, $\varphi(x)\in\overset{\circ}{H}{}^2(\Omega_0)$, а 


$f(x,t)\in L_2(Q_T)$ --- пространству функций с конечной нормой 


$$ 


\|u\|^2_{Q_T}=\int^T_0\int_{\Omega_0}u^2d\mu\,d t. 


$$ 


Обозначим через $\Omega_0 (t_1)$ сечение цилиндра $Q_T$ плоскостью 


$t=t_1$. 


Введем следующие функциональные пространства: 


\begin{itemize} 


\item[] $C^{i,1}_\sigma (Q_T)$ --- класс функций, принадлежащих 


$C^{i}_\sigma 


(\Omega_0 (t))$ по пространственным переменным для каждого $t \in 


[0,T]$ и гладких по $t$ в $\overline{Q}_T$, $i=1,2$; 


\item[] $\overset{\bullet}{C}{}^{i,1}(Q_T)$ --- подмножество функций из 


$C^{i,1}_\sigma (Q_T)\cap C(Q_T)$, равных нулю вблизи $S_T$; 


\item[] $\overset{\bullet}{C}{}^{2,1}_{rg}(Q_T)$ --- подмножество функций 


из $C^{2,1}_\sigma(Q_T)\cap C(Q_T)$, равных нулю вблизи $S_T$ и 


удовлетворяющих условию (8); 


\item[] $\overset{\circ}{H}{}^2(Q_T)$ --- 


пополнение $\overset{\bullet}{C}{}^{2,1}_{rg}(Q_T)$ 


по норме $\|\cdot\|_2^1$, порожденной скалярным произведением 


$$ 


\la u,v\ra_t=\int_{Q_T}{\bigg(\frac{\partial u}{\partial t}\frac{\partial v} 


{\partial t}+r\Delta_p u\Delta_p v+ 


q\nabla u \nabla v + k u v\bigg)}d\mu\,d t;


$$ 


\item[] $\overset{\bullet}{H}{}^2(Q_T)$ --- подмножество функций из 


$\overset{\circ}{H}{}^2(Q_T)$, обращающихся в нуль в обобщенном 


смысле при $t=T$. 


\end{itemize} 


Здесь $\nabla u$ и $\Delta_p u$ --- слабые градиент и лапласиан


соответственно, определяемые постратно. Функции, 


имеющие слабый градиент (лапласиан), 


будем называть слабо дифференцируемыми (дважды слабо дифференцируемыми). 


Аналогичным 


образом будем определять и слабую производную по времени. 





\noindent{}Слабым решением задачи (7)--(10) будем называть 


функцию $u(x,t) \in \overset{\circ}{H}{}^2(Q_T)$, равную 


$\varphi(x)$ при $t=0$ и удовлетворяющую (при условии, что $r,k,q\in 


L^2_\mu(\Omega_0)$, а $p$ --- слабо дифференцируемая на $\Omega_0$


функция) тождеству 


$$ 


\int_{Q_T}{(r\Delta_{p}u\Delta_{p}\eta+q\nabla 


u\nabla\eta+ku\eta-u_t \eta_t)}d\mu\,dt 


-\int_{\Omega_0}{\psi\eta(x,0)}d\mu=\int_{Q_T}{f\eta}\,d\mu\,dt 


$$ 


для любой функции 


$\eta \in \overset{\bullet}{H}{}^2(Q_T)$. 


 


Обозначим 


$$ 


a_k = \la \varphi,u_k\ra_\mu,\ \  


b_k = \frac{1}{\lambda_k}\la \psi,u_k\ra_\mu,\ \ 


T_k = \frac{1}{\lambda_k}\int_0^{t}\la f(x,\tau),u_k\ra_\mu 


\sin{\lambda_{k}(t-\tau)}d\tau. 


$$ 


Справедлива 


\vspace*{-1mm}


 


\begin{thm} 


Пусть $q(x)$ --- слабо дифференцируемая, а $r(x)$ ---


дважды слабо дифференцируемая 


на $\Omega_0$ функции. Тогда обобщенное решение 


задачи $(7)$--$(10)$ представимо в виде ряда 


$$ 


u(x,t)=\sum_{k=1}^{\infty}(a_k \cos{\lambda_{k}t}+ 


b_k\sin{\lambda_{k}t}+T_k (t))u_k (x), 


$$ 


сходящегося в $\overset{\circ}{H}{}^2(\Omega)$ равномерно по $t\in [0,T];$ 


его сумма принадлежит классу $\overset{\circ}{H}{}^2(Q_T)$. 


\end{thm} 


 


\begin{thebibliography}{99} 


\bibitem{1} 


Courant R. {\em \"Uber die Anwendung der Variationrechnung 


in der Theorie der Eigenschwingungen und \"uber neue Klassen von 


Funktionalgleichungen} // Acta Math. -- 1926. -- Bd.\,40. -- S.\,1--68. 


\bibitem{2} 


Ali Mehmeti F., von Below J., Nicaise S. (eds.) {\em Partial differential 


equations on multistructures}. -- Marcel Dekker, 2001. -- 258 p. 


\bibitem{3} 


Покорный Ю.В., Черникова Л.Н. {\em О краевой задаче на графе для $2$-звенной 


цепочки стержней} // Нелин. колеб. -- Ижевск, 1981. -- 


С.\,39--45. 


\bibitem{4} 


Покорный Ю.В., Пенкин О.М. {\em О теоремах сравнения для уравнений на 


графах} // Дифференц. уравнения. -- 1989. -- Т.\,25. -- \No\,7. -- 


С.\,1141--1150. 


\bibitem{5} 


Жиков В.В. {\em Связность и усреднение. Примеры фрактальной проводимости} 


// Матем. сб. -- 1996. -- Т.\,187. -- \No\,8. -- С.\,3--40. 


\bibitem{6} 


Жиков В.В. {\em Усреднение задач теории упругости на сингулярных 


структурах} // Изв. РАН. Сер. матем. -- 2002. -- Т.\,66. -- \No\,2. -- 


С.\,81--148. 


\bibitem{7} 


Пенкин О.М. {\em О принципе максимума для эллиптического уравнения на 


двумерном клеточном комплексе} // Докл. РАН. -- 1997. -- Т.\,352. -- 


\No\,4. -- С.\,462--465. 


\bibitem{8} 


Пенкин О.М., Покорный Ю.В. {\em О дифференциальных неравенствах для 


эллиптических уравнений на сложных многообразиях} // Докл. РАН. -- 1998. 


-- Т.\,360. -- \No\,4. -- С.\,456--458. 


\bibitem{9} 


Пенкин О.М., Покорный Ю.В. {\em О несовместных неравенствах для 


эллиптических уравнений на стратифицированных множествах} // Дифференц. 


уравнения. -- 1998. -- Т.\,34. -- \No\,8. -- С.\,1107--1113. 


\bibitem{10} 


Пенкин О.М., Богатов Е.М. {\em О слабой разрешимости задачи Дирихле на 


стратифицированных множествах} // Матем. заметки. -- 2000. -- Т.\,68. -- 


Вып.\,6. -- С.\,874--886. 


\bibitem{11} 


Богатов Е.М. {\em О разрешимости эллиптических уравнений на 


стратифицированных множествах}: Дис. $\dotsc$ канд. физ.-матем. наук. -- 


Воронеж, 2000. -- 91 с. 


\bibitem{12} 


Покорный Ю.В. {\em О знакорегулярных функциях Грина некоторых 


неклассических задач} // УМН. -- 1981. -- Т.\,36. -- Вып.\,4. -- 


С.\,205--206. 


\bibitem{13} 


Gavrilov A.A., Nicaise S., Penkin O.M. {\em Poincare's inequality on 


stratifies sets and applications}. -- Raport de recherche 01.2, Universite 


de Valencienes. -- P.\,1--20. 


\bibitem{14} 


Гаврилов А.А., Пенкин О.М. {\em Слабый принцип максимума для 


эллиптического оператора на стратифицированном множестве} // Тр. школы 


``Современные методы в теории краевых задач''. Ч.\,1. -- Воронеж, 2000. -- 


С.\,48--56. 


 


 


\end{thebibliography} 


 


\vskip 2cm 


 


\postup{\it Старооскольский технологический}{\qquad\it Поступили} 


\postup{\it институт $($филиал Московского }{\it первый вариант 


$15.02.2001$} 


\postup{\it института стали и сплавов$)$}{\it окончательный


вариант $10.07.2002$} 


\postup{\it Воронежский государственный университет}{} 


 


\label{end} 


\end{document} 


